% P01 manuscript skeleton - numbers are placeholders tied to artifact evidence
% IDs (E-xxx). Every numeric claim must be regenerated from
% research_artifacts_v3/runs/p01/ before submission (P01-T15/T17). No SOTA or
% learning-gain claim is permitted anywhere in this manuscript.
\documentclass[nonacm]{acmart}
\title{Risk-Controlled Selective Transfer Gates under Partial Cohort Alignment:
A Null Result on the Open University Learning Analytics Dataset}
\author{Anonymous Author(s)}
\affiliation{\institution{Anonymous Institution}\city{}\country{}}

\begin{abstract}
% OBSERVATIONAL STUDY. NO CAUSAL CLAIM. NO LEARNING-GAIN CLAIM.
We study whether a selective, risk-calibrated transfer gate avoids negative
transfer when a classifier trained on earlier presentations of the Open
University Learning Analytics Dataset (OULAD) is deployed on a held-out future
presentation (2014J) with only partially aligned cohorts. On 32{,}593
learner--module--presentation rows (Kuzilek et al., 2017; CC BY 4.0), a
regularized multinomial logistic baseline (C1) reaches an accuracy of
[E-C1-ACC] (95\% Wilson CI [E-C1-CI]) on the temporal-transfer test rows,
against [E-C0-ACC] for the majority sanity baseline (C0). Our risk-controlled
selective gate (RC-STG) with conformal-style thresholds calibrated only on a
calibration split achieves [E-PROP-ACC] at coverage [E-PROP-COV]; the
preregistered null---that the gate does not improve accuracy over the
strongest classical baseline---is NOT rejected (exact McNemar
$p=[E-MCN-P]$, Holm-adjusted; learner-cluster bootstrap 95\% CI on the paired
difference [E-BOOT-CI]). We report this as an explicit null result, with
ablations, calibration (ECE [E-ECE-PROP] vs.\ [E-ECE-C1]), worst-group
robustness across region, age band, disability, and deprivation decile, and a
coverage--risk accounting showing the gate's accuracy cost is concentrated in
its abstention behavior. Results are observational and single-dataset;
external generalization beyond one platform is unverified.
\end{abstract}

\keywords{knowledge tracing transfer, conformal risk control, selective
prediction, learning analytics, null result, OULAD}

\section{Introduction}
Knowledge tracing (KT) models predict a learner's evolving knowledge state
from interaction histories~\cite{corbett1994,piech2015}. Deploying a KT or
learner-outcome model trained on one platform or cohort onto another---where
only a subset of skills and behaviors align---risks \emph{negative transfer}.
The preregistered research question (P01-T01):
\emph{Can a selective risk-calibrated transfer gate avoid negative transfer
when the target platform has sparse learners and only partially aligned
skills?}

\section{Related Work}
\textbf{Learner-outcome datasets.} OULAD aggregates demographics, early VLE
clicks, assessment records and final outcomes for 32{,}593
learner--module--presentation registrations~\cite{kuzilek2017}; it supports
risk-ranking studies but its click logs are observational behavior proxies,
not measured self-regulation and not interventions.
\textbf{Knowledge tracing.} Bayesian knowledge tracing
\cite{corbett1994}; deep knowledge tracing \cite{piech2015}; attention-based
KT (SAKT/AKT) and cross-platform transfer models (e.g., DTransKT) target
interaction-level sequences; the required data shape does not exist in OULAD
(Section~\ref{sec:threats}).
\textbf{Selective prediction and risk control.} Conformal prediction
constructs distribution-free prediction sets under exchangeability
\cite{angelopoulos2021}; selective classification trades coverage for risk
\cite{geifman2017}; conformal risk control extends coverage guarantees to
bounded losses \cite{angelopoulos2022}. Our gate is a calibrated risk-control
\emph{proxy}: thresholds follow the conservative conformal quantile form, but
under presentation shift the exchangeability assumption only approximately
holds, so we measure the realized risk rather than claim a guarantee.

\section{Problem Setup}
Rows are learner--module--presentation records with a frozen 60-day early
click window and a 4-class outcome (Distinction/Fail/Pass/Withdrawn). The
label is sealed outside the feature vector; a leakage audit
(\texttt{research\_v3.eval.leakage}) must pass before any fitting. Features:
one-hot demographics and course fields over a train-pinned vocabulary plus
train-standardized numeric activity aggregates.

\subsection{Splits}
Two preregistered protocols (P01-T05), materialized once and hashed:
\textbf{T} (primary): chronological transfer, train on 2013B/2013J/2014B
(+ small 2014J tail), test on 2014J only; learners overlapping later
partitions are removed from earlier ones (no learner spans partitions).
\textbf{G} (secondary): learner-group split; all presentations appear in each
partition, no learner does.

\subsection{Models}
\textbf{C0 (sanity):} majority training label. \textbf{C1 (matched
classical):} $L_2$-regularized multinomial logistic regression, numpy,
full-batch gradient descent, hyperparameters frozen before the run
($\lambda{=}1.0$, 250 iterations, lr $0.5$; optimizer state
\texttt{heuristic}, no optimality certificate).
\textbf{RC-STG (proposed):} reuses C1 probabilities; an \textsc{Abstain}
state fires when the margin falls below a conservative conformal quantile of
correct-case calibration margins ($\alpha_{\text{conf}}{=}0.20$), and a Fail
re-route fires when $p(\text{Fail})$ exceeds a quantile of calibration
non-Fail rows ($\alpha_{\text{fail}}{=}0.30$), bounding the re-route
false-alarm rate at approximately $\alpha_{\text{fail}}$ on the calibration
distribution.

\section{Experimental Design}
Preregistered before the confirmatory run
(\texttt{research\_v3/experiments/p01/preregistration.yaml}); primary
endpoint: paired per-case accuracy on the frozen T-test rows; comparison
family \{proposed vs C1, proposed vs C0\}, Holm-corrected; exact McNemar
tests; learner-cluster bootstrap CIs (2{,}000 resamples); Wilson and
Clopper--Pearson intervals for all proportions. Ablations A1--A5 remove the
gate, its parts, or shuffle its thresholds (same split, same compute).

\section{Results}
% ALL NUMBERS BELOW ARE REGENERATED FROM research_artifacts_v3/runs/p01/.
\textbf{Primary (T, $n{=}[E-N-T]$).} C0 [E-C0-ACC], C1 [E-C1-ACC]
([E-C1-CI]), RC-STG [E-PROP-ACC] at coverage [E-PROP-COV]. The preregistered
null is not rejected; the gate is significantly \emph{worse} on raw accuracy
(McNemar $p=[E-MCN-P]$; bootstrap CI [E-BOOT-CI]) --- an explicit negative
result (P01-T14), disclosed per the preregistration.
\textbf{Mechanism check (A1--A5).} Removing the gate recovers C1 exactly
([E-A1-ACC]); degenerate gates (random/shuffled) degrade differently
([E-A4-ACC], [E-A5-ACC]), confirming the learned thresholds, not abstention
volume alone, drive the behavior.
\textbf{Calibration and robustness.} ECE: proposed [E-ECE-PROP], C1
[E-ECE-C1]. Worst-group accuracies (region/age/disability/IMD) for C1:
[E-WG-C1]; the gate's worst groups degrade with coverage [E-WG-PROP].
\textbf{Coverage--risk.} Among non-abstained rows the C1 backbone's accuracy
is [E-KEPT-ACC] vs.\ [E-C1-ACC] overall: the gate concentrates errors, but
the price is a [E-COV-LOSS] coverage loss that overwhelms the accuracy gain
in the preregistered per-case estimand, which counts abstentions as
incorrect for every system.

\section{Discussion}
The null supports the H0 disclosure obligation: selective risk calibration on
this task trades per-case accuracy for calibrated abstention without improving
the preregistered estimand. This is consistent with the coverage--risk
literature and is a useful negative result for risk-controlled learner-outcome
deployment: the accuracy cost of strict risk control under cohort shift can
exceed its benefit when the underlying model is only moderately discriminative.

\section{Threats to Validity}
\label{sec:threats}
\textbf{Observational, not causal.} OULAD is observational: early click
counts are behavior proxies, not interventions, and predicting
\texttt{final\_result} is not a learning gain. Nothing here estimates the
effect of any action on learners.
\textbf{Single dataset, single platform.} Only one external dataset
(\texttt{oulad-2013-2014}, APPROVED, CC BY 4.0) is available; spec-required
second-domain replication is BLOCKED and marked \emph{external generalization
unverified}.
\textbf{Comparator scope.} Deep KT baselines (DKT~\cite{piech2015}, SAKT,
AKT) and published cross-domain KT require interaction-level sequences
absent from the approved data shape; they are recorded as honest BLOCKED
cards (P01-T08), not silently substituted. C1 is therefore the strongest
\emph{executable} baseline, not the strongest published one.
\textbf{Shift breaks conformal assumptions.} Threshold semantics are
calibrated on source presentations; realized target risk (e.g., the Fail
false-alarm rate [E-FFA] vs.\ the nominal 0.30) quantifies the approximation.
\textbf{Human tasks.} Ethics/novelty/submission signoffs (P01-T19/T20/T21)
remain BLOCKED pending real human approval; no simulated signoffs exist.

\section{Conclusion}
A preregistered, leakage-audited, deterministic evaluation on real OULAD data
shows a conformal-style selective transfer gate does not beat a matched
regularized logistic baseline on the preregistered accuracy estimand. We
release all artifacts, seeds, and sealed-split manifests for independent
replication.

\begin{thebibliography}{9}
\bibitem{corbett1994} A.~T. Corbett and J.~R. Anderson, ``Knowledge tracing:
Modeling the acquisition of procedural knowledge,'' \emph{User Modeling and
User-Adapted Interaction}, 4(4):253--278, 1994.
\bibitem{kuzilek2017} J.~Kuzilek, M.~Hlosta, and Z.~Zdrahal, ``Open University
Learning Analytics dataset,'' \emph{Scientific Data}, 4:170171, 2017.
\bibitem{piech2015} C.~Piech et al., ``Deep knowledge tracing,'' \emph{NeurIPS
28}, 2015.
\bibitem{angelopoulos2021} A.~N. Angelopoulos and S.~Bates, ``A gentle
introduction to conformal prediction and distribution-free uncertainty
quantification,'' \emph{arXiv:2107.07511}, 2021.
\bibitem{angelopoulos2022} A.~N. Angelopoulos, S.~Bates, A.~Fisch, L.~Lei,
and T.~Schuster, ``Conformal risk control,'' \emph{arXiv:2208.02814}, 2022.
\bibitem{geifman2017} Y.~Geifman and R.~El-Yaniv, ``Selective classification
for deep neural networks,'' \emph{NeurIPS 30}, 2017.
\end{thebibliography}
\end{document}
